\subsection{Adapted covariance statistics, sketches, and tracking}
Gradients depend on previous updates, so their second moments form an adapted
process. We separate fluctuations around each conditional covariance from the
lag introduced by averaging a changing target. Normalized and raw EMAs are
distinguished throughout because their early-time scaling changes the relative
weights in the flag objective.

\begin{definition}[Filtration and normalized EMA]
For a filtration $(\mathcal F_t)$, let $G_t$ be $\mathcal F_t$-measurable and
\begin{gather}
 \EE_t[\cdot]:=\EE[\cdot\mid\mathcal F_{t-1}],\\
 X_t:=G_t^\top G_t,\quad C_t:=\EE[X_t\mid\mathcal F_{t-1}],\\
 Z_t:=X_t-C_t,\quad\EE[Z_t\mid\mathcal F_{t-1}]=0,\\
 \overline C_{k,t}:=(1-\rho_k)\sum_{i=1}^t\rho_k^{t-i}X_i,\quad
 \widehat C_{k,t}:=\overline C_{k,t}/(1-\rho_k^t),\quad0<\rho_k<1,\\
 \alpha_{k,t,i}:=\frac{(1-\rho_k)\rho_k^{t-i}}{1-\rho_k^t},\qquad
 C^\circ_{k,t}:=\sum_{i=1}^t\alpha_{k,t,i}C_i.
\end{gather}
\end{definition}
\begin{lemma}[Exact EMA identities]\label{lem:rigid-ema}
\begin{align}
 \sum_{i=1}^t\alpha_{k,t,i}&=1,\quad
 \max_i\alpha_{k,t,i}=\frac{1-\rho_k}{1-\rho_k^t},\\
 \sum_{i=1}^t\alpha_{k,t,i}^2
 &=\frac{1-\rho_k}{1+\rho_k}\frac{1+\rho_k^t}{1-\rho_k^t},\\
 A_{k,t}:=\sum_i(t-i)\alpha_{k,t,i}
 &=\frac{\rho_k}{1-\rho_k}-\frac{t\rho_k^t}{1-\rho_k^t},\\
 \|C_i-C_{i-1}\|_{\op}\le\nu
 &\Longrightarrow\|C^\circ_{k,t}-C_t\|_{\op}\le\nu A_{k,t}.
\end{align}
\end{lemma}
\begin{proof}
\begin{align}
 \sum_{j=0}^{t-1}\rho^j&=(1-\rho^t)/(1-\rho),\quad
 \sum_{j=0}^{t-1}\rho^{2j}=(1-\rho^{2t})/(1-\rho^2),\\
 \sum_{j=0}^{t-1}j\rho^j
 &=\rho\frac{\mathrm d}{\mathrm d\rho}\frac{1-\rho^t}{1-\rho},\\
 \|C^\circ_{k,t}-C_t\|_{\op}
 &\le\sum_i\alpha_{k,t,i}\sum_{j=i+1}^t\|C_j-C_{j-1}\|_{\op}
 \le\nu\sum_i(t-i)\alpha_{k,t,i}.
\end{align}
\end{proof}
\begin{theorem}[Finite-horizon adapted covariance bound]\label{thm:rigid-covariance}
Suppose, almost surely,
\begin{gather}
 \|Z_i\|_{\op}\le L_Z,\quad
 \EE[Z_i^2\mid\mathcal F_{i-1}]\preceq v_i I,
 \quad v_i\text{ deterministic},\\
 V_{k,t}:=\sum_i\alpha_{k,t,i}^2v_i,\quad
 R_{k,t}:=L_Z\max_i\alpha_{k,t,i},\\
 u:=\log(2dKT/\delta),\quad
 \mathfrak r(V,R,u):=\frac{Ru}{3}+\sqrt{2Vu+\frac{R^2u^2}{9}}.
\end{gather}
With probability at least $1-\delta$, simultaneously for $k\le K,t\le T$,
\begin{align}
 \|\widehat C_{k,t}-C_t\|_{\op}
 &\le\mathfrak r(V_{k,t},R_{k,t},u)+\nu A_{k,t},\\
 \|\overline C_{k,t}-C_t\|_{\op}
 &\le(1-\rho_k^t)
  \left(\mathfrak r(V_{k,t},R_{k,t},u)+\nu A_{k,t}\right)
       +\rho_k^t\|C_t\|_{\op}.
\end{align}
\end{theorem}
\begin{proof}
For fixed $(k,t)$, $D_i:=\alpha_{k,t,i}Z_i$ satisfies
\begin{gather}
 \EE[D_i\mid\mathcal F_{i-1}]=0,\quad
 \lambda_{\max}(\pm D_i)\le R_{k,t},\quad
 \left\|\sum_i\EE[D_i^2\mid\mathcal F_{i-1}]\right\|_{\op}\le V_{k,t}.
\end{gather}
The matrix Freedman inequality \citep{tropp2011freedman} gives
\begin{align}
 \PP\left(\left\|\sum_iD_i\right\|_{\op}\ge s\right)
 &\le2d\exp\left(-\frac{s^2}{2(V_{k,t}+R_{k,t}s/3)}\right),\\
 s=\mathfrak r(V_{k,t},R_{k,t},u)
 &\Longrightarrow s^2=2u(V_{k,t}+R_{k,t}s/3),\\
 \PP\left(\exists k,t:\|\widehat C_{k,t}-C^\circ_{k,t}\|_{\op}>s\right)
 &\le2dKT e^{-u}=\delta,\\
 \overline C_{k,t}-C_t
 &=(1-\rho_k^t)(\widehat C_{k,t}-C_t)-\rho_k^t C_t.
\end{align}
Thus the radius exactly inverts the Bernstein exponent in matrix Freedman;
it avoids the further relaxation
$\mathfrak r(V,R,u)\le\sqrt{2Vu}+2Ru/3$. A zero conditional-variance budget
means the martingale sum vanishes almost surely and its radius can be set to zero.
\end{proof}
\begin{corollary}[Memory--drift balance]
For the limiting EMA weights, $v_i\le v$, $h=1-\rho$, and fixed $u$,
\begin{gather}
 \epsilon(h)\le\mathfrak r\left(\frac{vh}{2-h},L_Zh,u\right)
                         +\frac{\nu(1-h)}h,\\
 \mathfrak r\left(\frac{vh}{2-h},L_Zh,u\right)
 =\sqrt{vu}\sqrt h+\frac{L_Zu}{3}h+O(h^{3/2}),\quad h\downarrow0,\\
 a:=\sqrt{vu}>0,\quad\nu>0,\quad
 h_0:=(2\nu/a)^{2/3}\le1,\quad
 a\sqrt{h_0}+\nu/h_0=3\,2^{-2/3}a^{2/3}\nu^{1/3}.
\end{gather}
The choice $h_0$ minimizes the leading stochastic--drift expression
$a\sqrt h+\nu/h$; the first display retains the finite-$h$ constants.
\end{corollary}
Longer memory reduces the stochastic term by spreading mass over more
observations, but increases the age of the information used to track the
current covariance. The displayed balance quantifies this tradeoff. Different
timescales can address different drift regimes; their joint benefit still
depends on the flag objective and the resulting update.

\begin{definition}[Fixed-probe Nystr\"om reconstruction]
\begin{gather}
 \Omega_k\in\RR^{d\times s_k},\quad s_k=\min(d,r_k+p),\quad
 Y_k=C_k\Omega_k,\quad B_k=\sym(\Omega_k^\top Y_k)+\gamma I,\quad\gamma>0,\\
 \mathcal N_\gamma(Y_k):=Y_kB_k^{-1}Y_k^\top.
\end{gather}
\end{definition}
\begin{proposition}[PSD order, regularization, and probe blindness]\label{prop:rigid-nystrom}
For $C\succeq0$, $Y=C\Omega$, and $Z=C^{1/2}\Omega$,
\begin{align}
 \mathcal N_\gamma(C\Omega)
 &=C^{1/2}Z(Z^\top Z+\gamma I)^{-1}Z^\top C^{1/2},\\
 0\preceq\mathcal N_\gamma(C\Omega)&\preceq C,\\
 \mathcal N_0(C\Omega)-\mathcal N_\gamma(C\Omega)
 &=C^{1/2}L\diag\left(\frac{\gamma}{\sigma_i^2+\gamma}\right)L^\top C^{1/2},
 \quad Z=L\diag(\sigma_i)R^\top\quad(\sigma_i>0).
\end{align}
If $s<d$, $v\in\ker\Omega^\top$, and $\|v\|=1$, then
\begin{equation}
 C=a vv^\top\quad\Longrightarrow\quad Y=0,\quad
 \mathcal N_\gamma(Y)=0,\quad\|C-\mathcal N_\gamma(Y)\|_{\op}=a.
\end{equation}
Here $\mathcal N_0$ uses the Moore--Penrose inverse.
\end{proposition}
\begin{proof}
\begin{align}
 Z(Z^\top Z+\gamma I)^{-1}Z^\top
 &=L\diag\left(\frac{\sigma_i^2}{\sigma_i^2+\gamma}\right)L^\top,\\
 0&\le\frac{\sigma_i^2}{\sigma_i^2+\gamma}\le1,\quad
 Z(Z^\top Z)^\dagger Z^\top=LL^\top,\\
 C\Omega&=av(v^\top\Omega)=0.
\end{align}
\end{proof}
\begin{lemma}[Sketch sensitivity]\label{lem:rigid-sketch-lip}
Suppose $B=\sym(\Omega^\top Y)+\gamma I\succeq\gamma I$ and
$B'=\sym(\Omega^\top Y')+\gamma I\succeq\gamma I$, with
$\max(\|Y\|_{\op},\|Y'\|_{\op})\le Y_\star$. Then
\begin{align}
 \|\mathcal N_\gamma(Y)-\mathcal N_\gamma(Y')\|_F
 &\le L_N\|Y-Y'\|_F,\quad
 L_N:=2Y_\star/\gamma+Y_\star^2\|\Omega\|_{\op}/\gamma^2,\\
 \mathrm D\mathcal N_\gamma(Y)[E]
 &=EB^{-1}Y^\top+YB^{-1}E^\top
                   -YB^{-1}\sym(\Omega^\top E)B^{-1}Y^\top.
\end{align}
\end{lemma}
\begin{proof}
\begin{align}
 B^{-1}-B'^{-1}&=B^{-1}(B'-B)B'^{-1},\\
 YB^{-1}Y^\top-Y'B'^{-1}Y'^\top
 &=(Y-Y')B^{-1}Y^\top+Y'B^{-1}(Y-Y')^\top
                       +Y'(B^{-1}-B'^{-1})Y'^\top,\\
 \|B^{-1}-B'^{-1}\|_F
 &\le\gamma^{-2}\|\Omega\|_{\op}\|Y-Y'\|_F,\qquad
 \mathrm D(B^{-1})[E]=-B^{-1}\sym(\Omega^\top E)B^{-1}.
\end{align}
\end{proof}
The probe-blindness example explains why sketch error must remain a separate
quantity: a fixed probe can miss an entire energetic direction. We now combine
this approximation error with statistical error and target motion. The
tracking recursion measures distance to a local oracle at refresh times, rather
than assuming exact covariance information at every parameter update.

\begin{definition}[Refresh-time chart recursion]
At refresh index $s$, let
\begin{gather}
 f_s(x):=\Phi(\psi(x);C_s),\quad \nabla f_s(x_s^\star)=0,\\
 x_{s+1}=x_s+\xi\nabla f_s(x_s)+\xi e_s+\xi\zeta_s+r_s,\\
 \Delta_s:=x_{s+1}^\star-x_s^\star,\quad
 E_s:=\EE\|x_s-x_s^\star\|^2,\\
 q:=\max\{|1-\xi\mu_\phi|,|1-\xi L_\phi|\}<1,\quad0<\xi<2/L_\phi.
\end{gather}
\end{definition}
\begin{theorem}[Cocoercive tracking with exact forcing gains]\label{thm:tf-tracking}
Under Assumption~\ref{ass:tf-chart}, assume the relevant iterates and connecting
segments remain in its convex chart. Define
\begin{gather}
 v_s:=\xi e_s+\xi\zeta_s+r_s-\Delta_s,\quad
 b_s:=\sqrt{\EE\|v_s\|^2}.
\end{gather}
Then the root-mean-square error obeys
\begin{align}
 \sqrt{E_{s+1}}&\le q\sqrt{E_s}+b_s,\\
 \sqrt{E_s}&\le q^s\sqrt{E_0}+\sum_{j<s}q^{s-1-j}b_j.
 \label{eq:tracking-rms-sharp}
\end{align}
The contraction is minimized at
$\xi_\star=2/(L_\phi+\mu_\phi)$, giving
$q_\star=(L_\phi-\mu_\phi)/(L_\phi+\mu_\phi)$.
For comparison, a squared-error form remains valid: for $\alpha>0$ with
$\kappa:=(1+\alpha)q^2<1$,
\begin{align}
 E_{s+1}&\le\kappa E_s+B_s,\\
 B_s&:=4(1+\alpha^{-1})\EE\left[
       \xi^2\|e_s\|^2+\xi^2\|\zeta_s\|^2+\|r_s\|^2+\|\Delta_s\|^2\right],\\
 E_s&\le\kappa^s E_0+\sum_{j=0}^{s-1}\kappa^{s-1-j}B_j,
 \label{eq:tf-tracking-expanded}\\
 \sum_{s=0}^{S-1}E_s&\le\frac{E_0+\sum_{s=0}^{S-2}B_s}{1-\kappa}.
\end{align}
If $\EE[\zeta_s\mid\mathcal H_s]=0$, where $\mathcal H_s$ contains
$(f_s,x_s,x_s^\star,e_s,r_s,\Delta_s)$, set
$d_s^2:=\EE\|\xi e_s+r_s-\Delta_s\|^2$ and
$\sigma_s^2:=\EE\|\zeta_s\|^2$. Then
\begin{equation}
 E_{s+1}\le(q\sqrt{E_s}+d_s)^2+\xi^2\sigma_s^2.
\end{equation}
For uniform envelopes $d_s\le\overline d$, $\sigma_s\le\overline\sigma$, define
\begin{equation}
 x_\infty:=\frac{q\overline d+
      \sqrt{\overline d^2+(1-q^2)\xi^2\overline\sigma^2}}{1-q^2}.
\end{equation}
Then $\sqrt{E_s}\le x_\infty+q^s(\sqrt{E_0}-x_\infty)_+$.
In particular the pure-drift RMS gain is $1/(1-q)$, whereas the centered-noise
variance gain is $\xi^2/(1-q^2)$; scalar quadratic recursions attain both.
\end{theorem}
\begin{proof}
The smooth strongly convex interpolation inequality associated with
Baillon--Haddad cocoercivity~\citep{bauschke2010baillon} gives the optimal-step
contraction. Here its local Hessian proof also gives every admissible step:
with $v=x-y$ and $g=\nabla(-f_s)(x)-\nabla(-f_s)(y)$,
\begin{align}
 \langle v,g\rangle&\ge\frac{\mu_\phi L_\phi}{\mu_\phi+L_\phi}\|v\|^2
                 +\frac{\|g\|^2}{\mu_\phi+L_\phi},\\
 T_s(x)-T_s(y)
 &=\left[I+\xi\int_0^1\nabla^2f_s(y+u(x-y))\,\mathrm du\right](x-y),\\
 \|T_s(x)-T_s(y)\|&\le q\|x-y\|,\quad T_s(x_s^\star)=x_s^\star,\\
 x_{s+1}-x_{s+1}^\star
 &=T_s(x_s)-T_s(x_s^\star)+\xi e_s+\xi\zeta_s+r_s-\Delta_s,\\
 \|a+b\|^2&\le(1+\alpha)\|a\|^2+(1+\alpha^{-1})\|b\|^2,\\
 \|b_1+b_2+b_3+b_4\|^2&\le4\sum_{i=1}^4\|b_i\|^2,\\
 \sum_{s=0}^{S-1}\sum_{j<s}\kappa^{s-1-j}B_j
 &\le(1-\kappa)^{-1}\sum_{j=0}^{S-2}B_j.
\end{align}
The averaged Hessian in the first map identity has spectrum in
$[-L_\phi,-\mu_\phi]$, so the endpoint formula for $q$ is exact for the
quadratic class. Minkowski's inequality in $L^2$ gives
\eqref{eq:tracking-rms-sharp} directly. Conditional centering removes the
cross term with $\zeta_s$. The scalar map
$x\mapsto\sqrt{(qx+\overline d)^2+\xi^2\overline\sigma^2}$ has Lipschitz constant at most
$q$ and fixed point $x_\infty$, proving the last bound. Constant scalar forcing
and a scalar AR(1) process respectively attain the two stated gains.
Sharpness here concerns the comparison recursion; applicability to a flag
trajectory still requires the stated chart invariance.
\end{proof}
\begin{corollary}[End-to-end flag forcing]\label{cor:rigid-flag-forcing}
Suppose $\|\widetilde C_{k,s}-\widehat C_{k,s}\|_{\op}\le a_{k,s}$ and
$\|\widehat C_{k,s}-C_{k,s}\|_{\op}\le\epsilon_{k,s}$. Then
\begin{align}
 \|e_s\|&\le2J_\psi\sum_k\omega_k\sqrt{r_k}(a_{k,s}+\epsilon_{k,s}),\\
 \|\Delta_s\|&\le\frac{2J_\psi}{\mu_\phi}
               \sum_k\omega_k\sqrt{r_k}\|C_{k,s+1}-C_{k,s}\|_{\op}.
\end{align}
If $\|C_k(w)-C_k(w')\|_{\op}\le L_{C,k}\|w-w'\|$ and
$w_{t+1}=w_t-\eta_t d_t$, with $\tau$ steps per refresh, then
\begin{align}
 \|\Delta_s\|^2
 &\le L_\star^2\tau\sum_{t=s\tau}^{(s+1)\tau-1}\eta_t^2\|d_t\|^2,
 \quad L_\star:=\frac{2J_\psi}{\mu_\phi}\sum_k\omega_k\sqrt{r_k}L_{C,k}.
\end{align}
\end{corollary}
\begin{proposition}[Bounded-gradient constants and raw-EMA weights]\label{prop:rigid-raw-ema}
If $\|G_t\|_F\le G_\star$ almost surely, then
\begin{gather}
 0\preceq X_t,C_t\preceq G_\star^2I,\quad
 \|Z_t\|_{\op}\le G_\star^2,\quad
 \EE_tZ_t^2=\EE_tX_t^2-C_t^2\preceq G_\star^4I,\\
 L_Z=G_\star^2,\qquad v_t=G_\star^4
 \quad\text{are admissible in Theorem~\ref{thm:rigid-covariance}},\\
 \Phi(U;\overline C_{1:K,t})
 =\sum_k\omega_k(1-\rho_k^t)\tr(P_k(U)\widehat C_{k,t}).
\end{gather}
Thus replacing raw EMAs by normalized EMAs changes the relative objective weights
unless all factors $1-\rho_k^t$ coincide.
\end{proposition}
\begin{proof}
\begin{gather}
 -G_\star^2I\preceq X_t-C_t\preceq G_\star^2I,\quad
 X_t^2\preceq G_\star^2X_t,\\
 \EE_t(X_t-C_t)^2=\EE_tX_t^2-C_t^2
 \preceq G_\star^2 C_t\preceq G_\star^4I.
\end{gather}
\end{proof}
The preceding chart recursion uses a Euclidean gradient for clarity. Actual
manifold ascent uses the inverse pullback metric, and QR introduces a retraction
remainder. The next theorem controls both effects so that the same tracking
argument applies under explicit local metric and step-size conditions.

\begin{theorem}[Actual Riemannian ascent in a non-Euclidean chart]
\label{thm:rigid-chart-metric}
Under Assumption~\ref{ass:tf-chart}, let $\mathcal D$ lie in the radius-$r$ ball
about $x^\star$, and define the pullback metric $\mathsf G(x)$ and natural gradient
\begin{gather}
 \mathsf G(x)_{ab}:=g_{\psi(x)}(\mathrm D\psi(x)e_a,\mathrm D\psi(x)e_b),\quad
 A(x):=\mathsf G(x)^{-1}\nabla f(x),\\
 \|\mathsf G^{-1}(x)-I\|_{\op}\le\delta_M,\quad
 \|\mathrm D(\mathsf G^{-1})(x)\|\le L_M,\\
 \mu_A:=\mu_\phi-L_\phi(\delta_M+L_Mr)>0,\quad
 L_A:=L_\phi(1+\delta_M+L_Mr),\\
 0<\xi<2\mu_A/L_A^2,\quad q_A:=\sqrt{1-2\mu_A\xi+L_A^2\xi^2}<1.
\end{gather}
Then
\begin{align}
 \langle x-y,A(x)-A(y)\rangle&\le-\mu_A\|x-y\|^2,\\
 \|A(x)-A(y)\|&\le L_A\|x-y\|,\\
 \|x+\xi A(x)-y-\xi A(y)\|&\le q_A\|x-y\|.
\end{align}
Let $\mathcal R_x(v)$ be the chart representation of the quotient QR retraction,
with $\mathcal R_x(0)=x$, $\mathrm D\mathcal R_x(0)=I$, and
$\|\mathrm D^2\mathcal R_x(v)\|\le C_\psi$ on the relevant segment. For an
estimated covariance $\widehat C$, the actual ascent is
\begin{gather}
 x^+=\mathcal R_x\left(\xi\mathsf G(x)^{-1}\nabla f_{\widehat C}(x)\right)
 =x+\xi A(x)+\xi e(x)+r(x),\\
 \|e(x)\|\le(1+\delta_M)E_C,\qquad
 \|r(x)\|\le\tfrac12 C_\psi\xi^2\|A(x)+e(x)\|^2.
\end{gather}
Theorem~\ref{thm:tf-tracking} therefore applies with $q$ replaced by $q_A$,
$\kappa=(1+\alpha)q_A^2<1$, and the displayed error and remainder bounds.
\end{theorem}
\begin{proof}
For $M(x):=\mathsf G(x)^{-1}$,
\begin{align}
 \mathrm DA(x)[v]&=\nabla^2f(x)v+(M(x)-I)\nabla^2f(x)v
                           +(\mathrm DM(x)[v])\nabla f(x),\\
 \|\nabla f(x)\|&=\|\nabla f(x)-\nabla f(x^\star)\|\le L_\phi r,\\
 v^\top\mathrm DA(x)v
 &\le-\mu_\phi\|v\|^2+\delta_M L_\phi\|v\|^2+L_ML_\phi r\|v\|^2
 =-\mu_A\|v\|^2,\\
 \|\mathrm DA(x)\|_{\op}&\le L_A,\\
 \|x-y+\xi(A(x)-A(y))\|^2
 &\le(1-2\mu_A\xi+L_A^2\xi^2)\|x-y\|^2,\\
 e(x)&:=M(x)(\nabla f_{\widehat C}(x)-\nabla f_C(x)),\\
 r(x)&=\int_0^1(1-s)\mathrm D^2\mathcal R_x(sv)[v,v]\,\mathrm ds,
 \quad v=\xi(A(x)+e(x)).
\end{align}
\end{proof}
